% ==STYLE-BEGIN==
\documentclass[UTF8,11pt]{ctexart}
\usepackage{amsmath,amssymb,amsthm}
\usepackage[a4paper,margin=2.4cm,top=2.3cm,bottom=2.3cm]{geometry}
\usepackage{fancyhdr}
\usepackage{enumitem}
\usepackage{hyperref}
\hypersetup{hidelinks}

% ---- 版式 ----
\linespread{1.15}
\setlength{\parskip}{0.35em}
\setlength{\parindent}{1.6em}
\setlist{nosep,leftmargin=2em}

% ---- 页眉 ----
\pagestyle{fancy}
\fancyhf{}
\fancyhead[L]{\small\coursename}
\fancyhead[R]{\small\thepage}
\renewcommand{\headrulewidth}{0.4pt}

% ---- 定理环境：中文用正体，标题加粗，按小节共用一条流水号 ----
\newtheoremstyle{cnthm}{0.7em}{0.7em}{\normalfont}{}{\bfseries}{}{0.5em}{}
\theoremstyle{cnthm}
\newtheorem{thm}{定理}[section]
\newtheorem{lem}[thm]{引理}
\newtheorem{prop}[thm]{命题}
\newtheorem{cor}[thm]{推论}
\newtheorem{defn}[thm]{定义}
\newtheorem{rem}[thm]{注}
\newtheorem{exa}[thm]{例}
\newtheorem{conj}[thm]{猜想}
\renewcommand{\proofname}{\bfseries 证明}
\providecommand{\coursename}{}
% ==STYLE-END==
\renewcommand{\coursename}{高等代数\,·\,1105}
\xeCJKDeclareCharClass{CJK}{"2150 -> "218F}
\title{高等代数 1105}
\date{}
\begin{document}
\maketitle

\section{等价关系}

目的: 给集合中的元素分类. 先要回答: 如何规定集合中两个元素"相同"?

\begin{defn}[关系]
设 $S$ 是一个非空集合, $S\times S=\{(a,b)\mid a,b\in S\}$ 的一个子集 $R$ 称为 $S$ 上的一个\textbf{关系}.
\end{defn}

\begin{exa}
$\mathbb{R}$ 上的"相等"关系 $=$ 可看作直线 $\{(x,x)\}$; "小于"关系 $<$ 可看作这条直线的上半平面. 这个关系既不自反, 也不对称.
\end{exa}

\begin{exa}
$R=\{(x,y)\in\mathbb{R}^2\mid |x-y|<1\}$ 可看作两条平行直线 $y=x\pm1$ 之间的带形区域. 这个关系有自反性、对称性, 但\textbf{没有传递性}.
\end{exa}

\begin{defn}[等价关系]
若 $S$ 上的关系 $R$ 满足

(1) 自反性: $\forall a\in S$, $(a,a)\in R$;

(2) 对称性: 若 $(a,b)\in R$, 则 $(b,a)\in R$;

(3) 传递性: 若 $(a,b)\in R$ 且 $(b,c)\in R$, 则 $(a,c)\in R$,

则称 $R$ 是 $S$ 上的\textbf{等价关系}.
\end{defn}

$\forall a\in S$, 记
\[
[a]=\{b\in S\mid (a,b)\in R\},
\]
称 $[a]$ 为 $a$ 的\textbf{等价类}.

\begin{prop}
若 $R$ 是 $S$ 上的等价关系, 则 $\forall a,b\in S$, 有 $[a]\cap[b]=\varnothing$ 或 $[a]=[b]$.
\end{prop}

\begin{proof}
设 $[a]\cap[b]\neq\varnothing$, 取 $c$ 属于这个交集, 则 $(a,c)\in R$, $(b,c)\in R$. 由对称性 $(c,b)\in R$, 再由传递性得 $(a,b)\in R$.

任取 $x\in[a]$, 则 $(a,x)\in R$; 由对称性 $(x,a)\in R$; 再由 $(a,b)\in R$ 与传递性得 $(x,b)\in R$, 即 $x\in[b]$. 故 $[a]\subseteq[b]$, 同理 $[b]\subseteq[a]$, 于是 $[a]=[b]$.
\end{proof}

上一讲的"线性等价"就是我们遇到的第一个等价关系.

\begin{exa}
有理数 $\mathbb{Q}=\Bigl\{\dfrac nm\ \Big|\ m,n\in\mathbb{Z},\ m\neq0\Bigr\}$, 规定
\[
\frac{n_1}{m_1}=\frac{n_2}{m_2}\iff n_1m_2=m_1n_2 .
\]
显然有自反性、对称性; 传递性由整数乘法的性质可得.
\end{exa}

\section{子空间}

\begin{defn}[子空间]
设 $W\subseteq K^m$. 若 $W$ 非空, 且对加法与数乘都封闭, 即
\[
\forall\alpha,\beta\in W:\ \alpha+\beta\in W;
\qquad
\forall k\in K,\ \forall\alpha\in W:\ k\alpha\in W ,
\]
则称 $W$ 是 $K^m$ 的一个\textbf{子空间}.
\end{defn}

记 $\alpha_1,\cdots,\alpha_s\in K^m$ 是一个向量组, 令
\[
L(\alpha_1,\cdots,\alpha_s)=\Bigl\{\sum_{i=1}^{s}k_i\alpha_i\ \Big|\ k_i\in K,\ i=1,\cdots,s\Bigr\} .
\]

\begin{thm}
$L(\alpha_1,\cdots,\alpha_s)$ 是一个子空间, 称其为\textbf{由 $\alpha_1,\cdots,\alpha_s$ 生成的子空间}.
\end{thm}

\begin{defn}[基与维数]
设 $W\subseteq K^m$ 是子空间, 向量组 $\alpha_1,\cdots,\alpha_s\in W$ 满足

(1) $\alpha_1,\cdots,\alpha_s$ 线性无关;

(2) $L(\alpha_1,\cdots,\alpha_s)=W$, 即 $\forall\beta\in W$, $\beta=k_1\alpha_1+\cdots+k_s\alpha_s$,

则称 $\alpha_1,\cdots,\alpha_s$ 是 $W$ 的一组\textbf{基}; $W$ 的任意一组基中向量的个数称为 $W$ 的\textbf{维数}.
\end{defn}

\begin{thm}
任何一个子空间 $W$ 都有基.
\end{thm}

\begin{proof}[筛选法]
若 $W=\{0\}$, 取空集为基(维数 $0$). 否则取非零向量 $\alpha_1\in W$.

若 $W=L(\alpha_1)$, 则 $\alpha_1$ 已是一组基. 否则取 $\alpha_2\in W\setminus L(\alpha_1)$, 则 $\alpha_1,\alpha_2$ 线性无关(否则 $\alpha_2\in L(\alpha_1)$, 矛盾).

一般地, 若已取到线性无关组 $\alpha_1,\cdots,\alpha_k$ 且 $L(\alpha_1,\cdots,\alpha_k)\neq W$, 取
\[
\alpha_{k+1}\in W\setminus L(\alpha_1,\cdots,\alpha_k),
\]
则 $\alpha_1,\cdots,\alpha_{k+1}$ 仍线性无关.

由于 $K^m$ 中线性无关的向量组最多含 $m$ 个向量, 这个过程必在有限步 $r\leq m$ 停止. 此时 $W=L(\alpha_1,\cdots,\alpha_r)$, 即 $\alpha_1,\cdots,\alpha_r$ 是 $W$ 的一组基.
\end{proof}

\end{document}
